% Source pins (SHA-256): PifoStatement.lean=fe9475ebfa2c462d27cb0a99781926074e1281396f2c9e8d32d6dfca788d814e; pifo-normalization.md=2d287de7924aa30479218da3e0a9a31577ecb97c5a34395c089f7abed98f61e7.
\section{Deciding equivalence}
\label{sec:deciding-equivalence}

The canonical characterization turns universal operational equivalence into a
finite structural comparison.  This section gives the extractor, accounts for
its cost, and then identifies an asymptotically optimal finite-support test
class.

\subsection{Structural extraction}
\label{sec:structural-decision}

An effective comparison is obtained by scanning two assigned paths.  Advance
while both paths choose the same child.  If they reach a common leaf, compare
their leaf ranks; at their first different child, compare their reference
ranks at that node.  The scan takes time linear in the path length and returns
\(\mathsf{eff}_S(x,y)\).

The rooted-triple extractor must do more than suppress a common unary prefix.
For pairwise-distinct \(x,y,z\), define
\(\mathsf{TripleDescriptor}(S;x,y,z)\) as follows.
\begin{enumerate}
\item Advance the three path cursors while all three select the same child.
      If they end in one common leaf, return a one-leaf flat descriptor with
      the three encountered leaf ranks.
\item At the first node where their selected children are not all equal,
      retain the three reference ranks at that node as synthetic root ranks.
\item For a \(1+1+1\) child split, route the three types to three synthetic
      singleton leaves.
\item For a \(2+1\) split, suppose \(u,v\) select the shared child.  Replace
      their entire residual subtrees by \emph{one} synthetic shared leaf.  Use
      leaf ranks
      \begin{equation}
        (r_u,r_v)=
        \begin{cases}
          (0,1),&\text{if }\mathsf{eff}_S(u,v)\text{ is }<,\\
          (0,0),&\text{if }\mathsf{eff}_S(u,v)\text{ is }=,\\
          (1,0),&\text{if }\mathsf{eff}_S(u,v)\text{ is }>.
        \end{cases}
        \label{eq:synthetic-pair-ranks}
      \end{equation}
      Route the remaining type to another singleton leaf; its leaf rank is
      immaterial.
\end{enumerate}
Evaluate \autoref{eq:rooted-obstruction} on this constant-size flat descriptor.
The result is \(R_S(x,y\mid z)\).  Unary transparency justifies the initial
scan, the two-type lemma justifies the synthetic shared leaf, and the
triple-flattening lemma identifies the descriptor's obstruction bit with that
of the original restriction.

\begin{example}[Why the residual pair must be collapsed]
\label{ex:residual-pair-collapse}
Take topology
\(\mathsf{node}[\mathsf{node}[\mathsf{leaf},\mathsf{leaf}],
\mathsf{leaf}]\) and paths
\begin{equation}
\begin{array}{c|cc|c}
 & \text{root }(\text{child},\text{rank})
 & \text{inner }(\text{child},\text{rank})
 & \text{leaf rank}\\
\hline
x&(0,0)&(0,2)&0\\
y&(0,2)&(1,0)&0\\
z&(1,1)&\text{--}&0
\end{array}
\label{eq:residual-pair-example}
\end{equation}
The first split is \(2+1\), but \(x\) and \(y\) end at different terminal
leaves.  Their residual effective comparison is \(>\), so the shared
synthetic leaf receives ranks \((1,0)\), while its flat root receives ranks
\((0,2,1)\) for \(x,y,z\).  Hence
\[
 \mathsf{eff}(x,z)=<,\qquad
 \mathsf{eff}(z,y)=<,\qquad
 \operatorname{comp}(<,<)=<\ne\mathsf{eff}(x,y),
\]
and \(R_S(x,y\mid z)\) is true.  Directly, \(F_S(xyz)=yzx\): the root removes
the reference inserted by \(x\), but the selected inner queue emits \(y\).
Looking only at terminal-leaf equality would incorrectly return false.
\end{example}

For valid input schedulers \(S,T\), the Boolean decision procedure streams
these local invariants.
\begin{enumerate}
\item For every unordered pair \(\{x,y\}\), compute and cache one orientation
      of \(\mathsf{eff}_S(x,y)\) and \(\mathsf{eff}_T(x,y)\); reject on the
      first mismatch.  The reverse orientation is determined automatically.
\item For every ordered pairwise-distinct triple \((x,y,z)\), compute
      \(R_S(x,y\mid z)\) and \(R_T(x,y\mid z)\) with the descriptor above and
      compare the bits immediately.  Reject on a mismatch; do not store the
      cubic tables.
\item Accept if both loops finish.
\end{enumerate}
For \(k=0\) or \(k=1\), the relevant loops are empty, as required by the base
cases of the canonical form.

\begin{theorem}[Structural decision]
\label{thm:structural-decision}
For valid schedulers \(S,T\), the procedure accepts exactly the
interleaved-equivalent pairs:
\begin{equation}
 \mathsf{Decide}(S,T)=\mathsf{true}
 \quad\Longleftrightarrow\quad
 S\intereq T.
 \label{eq:structural-decision-correct}
\end{equation}
\end{theorem}

\begin{proof}
Acceptance says exactly that the effective-comparison tables and all rooted
obstruction bits agree.  Apply \autoref{thm:canonical-characterization}.  The
two finite loops also prove termination.
\end{proof}

\begin{theorem}[Decision complexity]
\label{thm:decision-complexity}
For valid input schedulers \(S,T\), let \(h\) be the larger of \(1\) and the
maximum number of constructors in an assigned path of either input, including
its terminal leaf; thus \(h=1\) when \(k=0\).  On a unit-cost RAM, structural
equivalence is decidable in
\begin{equation}
  O(k^3h)\ \text{time}
  \qquad\text{and}\qquad
  O(k^2)\ \text{working space}
  \label{eq:decision-complexity}
\end{equation}
beyond the immutable input.
\end{theorem}

\begin{proof}
There are \(\Theta(k^2)\) pair scans, each of length at most \(h\), and their
ternary results occupy \(O(k^2)\) entries.  There are
\(\Theta(k^3)\) triple scans.  A common-prefix scan, together with at most one
pair-suffix scan if the cache were absent, visits \(O(h)\) path cells; with the
cache, \autoref{eq:synthetic-pair-ranks} is constant time.  Each pair of
obstruction bits is discarded after comparison, so no cubic table is stored.
An iterative cursor implementation uses constant transient space; a recursive
one adds an \(O(h)\) call stack.
\end{proof}

If ranks and child indices have maximum bit width \(B\), charging comparison
by bit complexity gives the conservative time bound \(O(k^3hB)\).  The bound
in \autoref{thm:decision-complexity} is for the Boolean decision only: it does
not hide construction of \(\mathcal N\) or synthesis of concrete ranks.  For
comparison, if all \(R\)-bits are first stored, a direct decomposition-tree
construction scans
\begin{equation}
  O\!\left(k^3h+\sum_X |X|^3\right)
  \subseteq O(k^3h+k^4)
  \label{eq:normal-form-construction-cost}
\end{equation}
over the laminar family of canonical blocks \(X\), using
\(\Theta(k^3)\) extra storage.  There are at most \(k-1\) non-singleton
blocks, which gives the displayed \(O(k^4)\) summand.  Re-extracting rather
than storing the bits gives the looser \(O(k^4h)\) bound.  No separate
complexity claim is made for choosing integer witnesses to cell feasibility.

\subsection{Two- and three-push tests suffice}

Define the finite suite
\begin{equation}
 \mathcal T(A)=
 \{\flushOps(w): w\text{ is injective and }2\le |w|\le3\}.
 \label{eq:three-push-suite}
\end{equation}
It has
\begin{equation}
 k(k-1)+k(k-1)(k-2)=\Theta(k^3)
 \label{eq:three-push-suite-size}
\end{equation}
tests, each with at most three pushes and three pops.  Singleton drains are
fixed by packet conservation and supply no test information.

\begin{theorem}[Two- and three-push sufficiency]
\label{thm:three-push-sufficiency}
Suppose valid \(S,T\) agree on every test in \(\mathcal T(A)\).  Then
\(S\intereq T\).
\end{theorem}

\begin{proof}
Fix distinct \(x,y\).  The tests \(xy\) and \(yx\), after restricting and
consistently relabelling the two types, recover
\(\mathsf{eff}(x,y)\) by \autoref{tab:binary-recovery}.  Thus the two effective
tables agree.

Now fix distinct \(x,y,z\).  Restrict and relabel these types as \(\Fin(3)\).
The triple-flattening lemma supplies a valid flat representative for each
restriction; this is a semantic replacement, not an assertion that the
original restrictions are flat.  Their length-two and length-three injective
drains agree by hypothesis, and their singleton drains agree automatically.
The finite three-label recovery lemma says that these fifteen nonempty
injective drains determine the effective table and all labelled obstruction
bits of a flat realization.  Transporting back gives
\(R_S(x,y\mid z)=R_T(x,y\mid z)\).  Varying the pair and triple gives equal
\((\mathsf{eff},R)\), so \autoref{thm:canonical-characterization} applies.  When
\(k<2\), the empty or singleton canonical base case gives the conclusion.
\end{proof}

Thus repetitions are unnecessary: one occurrence of each of at most three
types recovers the complete invariant.

\subsection{A cubic support lower bound}

Pair data alone is insufficient.  The following explicit family will also
force every unordered triple support to appear in any universal suite whose
tests use at most three types.

\begin{example}[Equal pair behavior, different triple behavior]
\label{ex:cubic-support-witness}
On types \(x,y,z\), use the flat topology
\(\mathsf{node}[\mathsf{leaf},\mathsf{leaf},\mathsf{leaf}]\).  In type order
\(x,y,z\), let \(\ell\), \(\rho\), and \(\lambda\) denote the leaf, root-rank,
and leaf-rank vectors:
\begin{equation}
\begin{array}{c|ccc}
 & \ell & \rho & \lambda\\
\hline
 S  &(0,0,1)&(0,2,1)&(1,0,0)\\
 S' &(0,1,1)&(1,0,2)&(0,1,0)
\end{array}
\label{eq:cubic-support-witness}
\end{equation}
The third leaf is initially unused.  Both schedulers have
\begin{equation}
 \bigl(\mathsf{eff}(x,y),\mathsf{eff}(x,z),
       \mathsf{eff}(z,y)\bigr)=(>,<,<).
 \label{eq:cubic-support-effective}
\end{equation}
Consequently their restrictions to every pair agree on all operation words.
However, \(R_S(x,y\mid z)\) is true and
\(R_{S'}(z,y\mid x)\) is true, giving component partitions
\(\{\{x,y\},\{z\}\}\) and \(\{\{x\},\{y,z\}\}\), respectively.  The
distinction is visible in one flush:
\begin{equation}
  F_S(xyz)=yzx,
  \qquad
  F_{S'}(xyz)=zxy.
  \label{eq:cubic-support-distinguishing-flush}
\end{equation}
For example, \(S\)'s root first selects the reference inserted by \(x\), but
the shared child emits \(y\); this is the same reference-hijacking phenomenon
as in \exref{ex:residual-pair-collapse}.
\end{example}

\begin{theorem}[Support-level optimality]
\label{thm:support-lower-bound}
For test suites whose individual tests use at most three packet types, every
universal equivalence suite over \(k\) types needs at least
\begin{equation}
  \binom{k}{3}=\Omega(k^3)
  \label{eq:support-lower-bound}
\end{equation}
distinct triple supports.  Together with \autoref{eq:three-push-suite}, the
support complexity is \(\Theta(k^3)\).
\end{theorem}

\begin{proof}
Embed \exref{ex:cubic-support-witness} on an arbitrary target triple
\(\{a,b,c\}\).  Route every outsider identically in both schedulers to the
common unused third leaf, give it root rank \(3\), and give it leaf rank \(0\);
retain target ranks \(0,1,2\).  The choice \(3\) is load-bearing.  Every
target--outsider effective comparison is identical, and an outsider cannot
obstruct a shared target pair because the two comparisons through it have
opposite strict signs.  A target likewise cannot obstruct a pair of
outsiders.

Therefore the two extended schedulers agree on every restriction omitting at
least one of \(a,b,c\), while their target-triple flush still differs.  A test
using at most three types can cover only one unordered triple, so a universal
suite must cover every one of the \(\binom{k}{3}\) choices.
\end{proof}

This is a lower bound on supports.  It does not say that a particular order of
each triple is sufficient, that every pair word is indispensable, or that
\(\mathcal T(A)\) is inclusion-minimal.

\subsection{Orientations and shortened drains}

The six orders of one triple cannot be replaced by a fixed choice of three.
In fact, any selected orientation can be the only distinguishing one.

\begin{example}[A uniquely distinguishing triple orientation]
\label{ex:missing-orientation}
Consider the following flat realizations, with all vectors in type order
\(0,1,2\):
\begin{equation}
\begin{array}{c|ccc}
 & \text{leaf} & \text{root rank} & \text{leaf rank}\\
\hline
 U&(0,0,2)&(2,1,1)&(2,2,2)\\
 V&(0,1,0)&(0,0,0)&(1,0,0)
\end{array}
\label{eq:missing-orientation-pair}
\end{equation}
Both have effective table
\((\mathsf{eff}_{01},\mathsf{eff}_{02},\mathsf{eff}_{12})=(=,>,=)\), while
their obstruction-edge tuples \((01,02,12)\) are
\((\mathsf{true},\mathsf{false},\mathsf{false})\) and
\((\mathsf{false},\mathsf{true},\mathsf{false})\).  All ordered-pair drains
agree:
\begin{equation}
  01\mapsto01,\quad 02\mapsto20,\quad 10\mapsto10,\quad
  12\mapsto12,\quad 20\mapsto20,\quad 21\mapsto21.
  \label{eq:missing-orientation-pairs}
\end{equation}
Five of the six triple drains also agree; exactly \(012\) differs:
\begin{equation}
\begin{array}{c|cc}
\hline
\text{push order}&F_U&F_V\\
\hline
012&021&210\\
021&201&201\\
102&120&120\\
120&120&120\\
201&201&201\\
210&210&210\\
\hline
\end{array}
\label{eq:missing-orientation-drains}
\end{equation}
\end{example}

\begin{proposition}[No three-orientation subfamily]
\label{prop:no-three-orientations}
For each of the six orders of three fixed types, there are valid flat
schedulers that differ on exactly that triple drain while agreeing on the
other five and on all ordered-pair drains.  Hence no fixed choice of three
triple orientations, even together with all pair drains, is universally
sufficient.
\end{proposition}

\begin{proof}
Relabelling types by a permutation relabels every successful output, by an
induction on \(\run\).  It can therefore move the sole differing word \(012\)
in \exref{ex:missing-orientation} to any chosen orientation while preserving
agreement on the other five and on every pair observation.
\end{proof}

Thus each orientation is individually indispensable when the other five and
all pair observations are retained.  This is a statement about the
triple-orientation subfamily, not inclusion-minimality of a mixed global
suite.

One observation in every full drain is nevertheless redundant.

\begin{proposition}[One-pop redundancy]
\label{prop:one-pop-redundancy}
Let \(w\) be an injective push word of length \(n>0\).  For valid schedulers,
agreement after pushing \(w\) and performing \(n-1\) pops is equivalent to
agreement on the full \(n\)-pop flush.  For \(n=1\), no observed pop is needed.
\end{proposition}

\begin{proof}
Every valid push adds one leaf packet and one matching reference at each
internal node on its path.  The balance invariant makes all \(n\) drain pops
succeed, with outputs a permutation of the known input types.  After the
first \(n-1\) outputs, the last output is the unique input type not yet
emitted.  Thus the prefix determines the full drain, and equality of prefixes
is equivalent to equality of full drains.  When \(n=1\), the sole output is
already determined by the input.
\end{proof}

Consequently every test in \(\mathcal T(A)\) may omit its final pop without
losing information.  No claim is made about the globally sparsest mixture of
orientations, pair tests, and still shorter prefixes.
